\documentclass[10pt,a4paper,openany]{article}


\usepackage[english,italian]{babel} %lingue utilizzabili

%PER GRAFICA
\usepackage{graphicx}
\usepackage{fancyhdr} %per intestaz e piè di pagina
\usepackage{subfig}
\usepackage{float}


\usepackage[utf8]{inputenc} %codifica e input

\usepackage{hyperref} %indice linkato ai paragrafi di riferimento
\hypersetup{
  colorlinks,
  citecolor=gray,
  filecolor=brown,
  linkcolor=black,
  urlcolor=cyan
}

\usepackage{lastpage} %per sapere il numero di pagine totali

\usepackage{parcolumns} %testo affiancato in colonna

\usepackage{listings} %Per inserire codice
\usepackage[usenames]{color} %Per permettere la colorazione dei caratteri 

\usepackage[a4paper,top=3cm,bottom=2cm,left=2cm,right=2cm] {geometry} %imposta pagina formato A4 e setta margini a piacimento
\usepackage[pdftex]{lscape} %per settare pagine in landscape

\renewcommand{\familydefault}{\sfdefault}   %mette tutto il carattere con uno stile ''grazioso''

\definecolor{editorGray}{rgb}{0.95, 0.95, 0.95}
\definecolor{editorOrange}{rgb}{1, 0.5, 0} % #FF7F00 -> rgb(239, 169, 0)
\definecolor{editorGreen}{rgb}{0, 0.6, 0} % #007C00 -> rgb(0, 124, 0)




%personalizzazione classe per il codice XML
\lstnewenvironment{xml}[1][] 
{\lstset{basicstyle=\scriptstyle \ttfamily, columns=fullflexible, keywordstyle=\color{blue}\bfseries, ndkeywordstyle=\color{blue}\bfseries , ndkeywords={references}, numberstyle=\color{red}, commentstyle=\color{editorGreen}, showstringspaces=false, stringstyle=\color{editorOrange},
language=XML, basicstyle=\small,
numbers=left, numberstyle=\tiny,
tabsize=2, stepnumber=10, numbersep=5pt, breaklines=true, frame=single, rulecolor=\color{black}, #1}}
{\lstset {language=XML,morekeywords={xml}}}{}


%personalizzazione classe per il codice HTML
\lstnewenvironment{html}[1][] 
{\lstset{basicstyle=\scriptstyle \ttfamily, columns=fullflexible,
keywordstyle=\color{blue}\bfseries, ndkeywords={content,=,charset=, id=, width=, height=},	 ndkeywordstyle=\color{editorGreen}\bfseries , numberstyle=\color{red} commentstyle=\color{red}, showstringspaces=false, stringstyle=\color{editorOrange},
language=HTML, basicstyle=\small,
numbers=left, numberstyle=\tiny,
tabsize=2, stepnumber=10, numbersep=5pt, breaklines=true, frame=single, rulecolor=\color{black}, #1}}{}

%personalizzazione classe per il codice PHP
\lstnewenvironment{php}[1][] 
{\lstset{basicstyle=\scriptstyle \ttfamily, columns=fullflexible,
keywordstyle=\color{blue}\bfseries, ndkeywords={content,=,charset=, id=, width=, height=},	 ndkeywordstyle=\color{editorGreen}\bfseries , numberstyle=\color{red} commentstyle=\color{red}, showstringspaces=false, stringstyle=\color{editorGreen},
language=PHP, basicstyle=\small,
numbers=left, numberstyle=\tiny,
tabsize=2, stepnumber=10, numbersep=5pt, breaklines=true, frame=single, rulecolor=\color{black}, #1}}{}





\begin{document}

\begin{figure}
\centering
\includegraphics[angle=0,scale=.30]{LOGO.png}
\end{figure} 


\title{ \textbf{\huge PROGRAMMAZIONE AD OGGETTI}\vspace{0.7cm} \\ {\Huge Relazione del progetto} \vspace{0.3cm} \\ {\Large Anno accademico 2015/2016} }
%\date{}
 
\author{\textbf{Zecchin Giacomo (1070122)}}

\maketitle
\thispagestyle{empty}


%da qui stile intestazione e piè di pagina cosi definiti
\pagestyle{fancy}
\lhead{}
\rfoot{\vspace {1pt} Relazione del progetto}
\cfoot{\vspace{1pt} Pagina: \thepage\ di \pageref{LastPage}}
\fancyfoot[L]
	{\vspace{1pt}
   		\hyperref[sec:index] {\includegraphics[scale=0.05]{leftfoot.png}} %icona linkata all'index
	 }
\renewcommand{\footrulewidth}{0.5pt}

%%%%%%%%%%%%%%%%%%%%%%%%%

\tableofcontents \label{sec:index} %CREA INDICE LINKABILE



\newpage

\section{Abstract}


%\selectlanguage{italian}	%%%%inserirlo causerebbe errore in compilazione per lettere accentate-->non risolto


	
%per inserire un'immagine
%\begin{landscape}
%\begin{figure}
%\subsection{Schema E-R iniziale del database}
%\centering
%\includegraphics[angle=0,scale=.40]{projBasi.jpg}
%\hspace{1in}
%\label{schema}
%\caption{schema E-R}
%\end{figure}
%\end{landscape}










\end{document}
